\section{Method}\label{sec:method}
In this paper, we focus on VidQA, which can be generalized to other video understanding tasks such as video summarization and grounding~\cite{xiao2024can,liu2025timecraft}.
This task is formulated as \((V, q) \mapsto \text{answer}\), where \(V=\{f_i\}_{i=1}^{n}\) is a video with \(n\) frames, \(f_i\) is the \(i\)-th frame, and \(q\) is the query.
The index set of frames is denoted by \(\N_n := \{1, 2, \dots, n\}\), 
and the list-wise $k$ frame selection can be expressed as: 
\begin{equation}\label{eq:score}
    S^* = \mathop{\mathrm{argmax}}_{S \subseteq \N_n , \left|S\right|=k}\ \mathrm{Score}(S)\,.
\end{equation}
This problem presents two challenges:
1) How can we design an effective \(\mathrm{Score}\) function that simultaneously addresses query relevance, list-wise diversity, and sequentiality?
2) How can we tackle the NP-hard challenges of the list-wise subset frame selection problem to find an effective solution with practical time complexity?


In \cref{subsec:method_sim}, we introduce the conditional multiple Gaussian kernel~(CMGK) within the reproducing kernel Hilbert space~(RKHS) to estimate high-dimensional similarities between frames conditioned on the current query.
In \cref{subsec:method_dpp}, we employ the determinantal point process~(DPP) based on the conditional similarity matrix of the frames, capturing both query relevance and list-wise frame diversity. 
Moreover, in \cref{subsec:method_kdp}, we partition the video into continuous segments and model frame selection as a Markov decision process~(MDP) to ensure list-wise diversity within local segments while maintaining sequentiality across segments.



\subsection{High-dimensional Conditional Similarity}\label{subsec:method_sim}
With the success of VLMs~\cite{radford2021learning,zhai2023sigmoid,alayrac2022flamingo},
frames and queries can be mapped into a unified latent embedding space.
For conciseness, we reuse \(f_i\) and \(q\) to denote the embeddings of the \(i\)-th frame and the query, respectively, to avoid defining redundant symbols. 
This is expressed as:
\begin{equation}\label{eq:emb} 
	f_i \leftarrow \mathrm{VLM_{vision}}(f_i) \in \mathbb{R}^d, \ q \leftarrow \mathrm{VLM_{text}}(q) \in \mathbb{R}^d \,,
\end{equation}
where \(\mathrm{VLM_{vision}}()\) and \(\mathrm{VLM_{text}}()\) denote the vision and text encoders of the VLM, respectively.
Notably, since the vision encoder is commonly used in MLLMs, the additional parameters introduced by the text encoder are negligible.


Instead of directly estimating similarity using cosine similarity or the inner product of embeddings, we map the embeddings into an RKHS via the reproducing kernel feature map \(\phi \in \mathcal{H}_k\), 
where, according to the kernel trick~\cite{scholkopf2005support}, the high-dimensional similarity in RKHS is expressed by \(\left<\phi(x), \phi(y)\right> = k(x, y)\).
%
We employ a characteristic multi-kernel defined as a convex combination of multiple positive semi-definite~(PSD) kernels \(\{k_u\}\)~(\eg Gaussian kernel): 
\begin{equation}\label{eq:k_family}
	\mathcal{K} \triangleq \left\{ 
		k = \sum_{u=1}^{U} \beta_u k_u : \sum_{u=1}^{U} \beta_u = 1, \beta_u \geq 0, \forall u 
	\right\}\,.
\end{equation}
The multi-kernel utilizes various kernels to enhance the estimation of high-dimensional similarity, thereby providing a principled method for optimal kernel selection~\cite{long2015learning,sun2023enhancing}.


In VidQA tasks, both frame-query relevance and list-wise diversity among the selected frames are important. 
Therefore, estimating frame similarity should consider the query relevance.
To address this, we propose the CMGK as:
\begin{equation}\label{eq:cond_k}
    \tilde{k}(f_i, f_j | q) = g(f_i, q)k(f_i, f_j)g(f_j, q),\quad g, k \in \mathcal{K}\,.
\end{equation}
This conditional kernel refines frame similarity by considering the query, reducing the importance of irrelevant frames.


Consequently, we construct a high-dimensional conditional similarity matrix that captures the similarities among frames while considering the query relevance, expressed as:
\begin{equation}\label{eq:tilde_L}
    \tilde{\mL} \in \mathbb{R}^{n \times n}, \quad \tilde{L}_{ij} = \tilde{k}(f_i, f_j | q)\,.
\end{equation}
This can also be expressed as:
\begin{equation}\label{eq:cmgk}
    \tilde{\mL} = \mathrm{diag}(\mr)\cdot \mL \cdot \mathrm{diag}(\mr)\,,
\end{equation}
with frame similarities \(\mL \in \mathbb{R}^{n \times n}, \mL_{ij} = k(f_i, f_j), k \in \mathcal{K}\), and frame-query relevances \(\mr \in \mathbb{R}^n, \mr_i = g(f_i, q), g \in \mathcal{K}\).


\subsection{Determinantal Point Process~(DPP)}\label{subsec:method_dpp}

DPP was initially introduced to model fermion repulsion in quantum physics and has proven effective in modeling list-wise diversity in machine learning~\cite{ye2023compositional,chen2018fast}.
We apply DPP to the high-dimensional conditional similarity matrix \(\tilde{\mL}\) to capture both query relevance and list-wise diversity in frame selection.
In fact, it represents the geometric volume of the subset in the RKHS; a larger volume indicates a more well-spread subset, leading to better list-wise diversity. 


Formally, DPP \(\mathcal{P}\) serves as a probability measure over the \(2^n\) subsets of frames sampled without replacement from a video \(V = \{f_i\}_{i=1}^{n}\).
The probability of selecting a subset of frame indices \(S \subseteq \N_n\) is defined as:
\begin{equation}\label{eq:dpp}
	\mathcal{P}(S)=\frac{\det(\tilde{\mL}_S)}{\det({\tilde{\mL} + \I})} \propto \det(\tilde{\mL}_S) \,,
\end{equation}
where \(\tilde{\mL}_S \equiv [\tilde{\mL}_{ij}]_{i,j \in S} \in \mathbb{R}^{|S|\times|S|}\) is a PSD matrix denoting the submatrix of \(\tilde{\mL}\) corresponding to the selected frame indices in \(S\).
%
Based on \cref{eq:cmgk} and \cref{eq:dpp}, we can derive the unnormalized log-probability of the subset \(S\) as:
\begin{equation}
	\log\left(\det(\tilde{\mL}_S)\right) = \sum_{i\in S}\log\left(\mr_i^2\right) + \log\left(\det(\mL_S)\right)\,,
\end{equation}
which clearly shows how to capture query relevance~(\ie \(\mr\)) and list-wise frame diversity~(\ie \(\det(\tilde{\mL}_S)\)) in frame selection.
Without loss of generality, we introduce a weighting factor to trade-off query relevance and list-wise diversity:
\begin{equation}\label{eq:log_dpp}
	\log \left( \det(\tilde{\mL}_S) \right) = \frac{1}{\lambda} \sum_{i \in S} \log\left(\mr_i^2\right) + \log\left(\det(\mL_S)\right)\,.
\end{equation}
This is a more general form of the standard DPP, equivalent to scaling the bandwidth \(h\) in the Gaussian kernel \mbox{\(k(x,y) = \exp\left(-\frac{\| (x-y)/h \|^2}{2\sigma^2}\right)\)} from \(h\) to \(\sqrt{\lambda} \cdot h\).



Frame selection follows the law of diminishing marginal utility, 
as the marginal gain from selecting additional frames is non-increasing~(as shown in \cref{fig:gain_on_videomme}).
This phenomenon is formally known in statistics as \emph{submodularity}, 
which is a well-established property of DPP, as shown in previous research~\cite{fujishige2005submodular,krause2014submodular}.
%
Selecting the optimal subset to maximize the submodular score function~\cref{eq:dpp,eq:log_dpp}, known as the submodular maximization problem in subset selection, is generally NP-hard.
Fortunately, there is a popular greedy algorithm can solve this problem,
which offers polynomial complexity and guarantees a \((1 - 1/e)\)-approximation, \ie, the solution is at least \((1 - 1/e)\) of the optimal score~\cite{chen2018fast}.

Specifically, we employ maximum a posteriori~(MAP) inference~\cite{chen2018fast,ye2023compositional} to solve DPP~(corresponding to \cref{eq:dpp,eq:log_dpp}).
We initialize the selected subset as \(S \leftarrow \emptyset\).
In each iteration, we select a new frame that maximizes the marginal gain in log-probability, given by:
\begin{equation}\label{eq:greedy}
	j = \mathop{\mathrm{argmax}}\limits_{i\in \mathbb{N}_n \backslash S} \log\left(\det(\tilde{L}_{S \cup \{i\}})\right) - \log\left(\det(\tilde{L}_S)\right)\,.
\end{equation}
The selected frame index \(j\) is then added to \(S\) until the subset size reaches the capacity limit \(k\).
%
By using Cholesky decomposition, the time complexity is reduced from \(\mathcal{O}(nk^3)\) to \(\mathcal{O}(nk)\) per iteration through incremental updates to the Cholesky factor~\cite{chen2018fast}.
Therefore, the overall time complexity of the MAP inference is \(\mathcal{O}(nk^2)\), compared to the vanilla KNN, the additional latency is negligible while \(k \ll n\)~\cite{ye2023compositional}.


\subsection{Markov Decision DPP}\label{subsec:method_kdp}

As described above, applying DPP to the conditional similarity matrix effectively captures both query relevance and list-wise diversity.
However, standard DPP treats video frames as independently permutable, disregarding the sequential structure inherent in video.
%
To address this, existing works on video summarization~\cite{gong2014diverse,zheng2021k} segment a video into multiple consecutive short segments and apply DPP within each segment, conditioned on the selections from the preceding segment.
This approach effectively models sequentiality by \emph{enhancing diversity among neighboring frames while reducing it for temporally distant ones.}
%
In VidQA, frames relevant to the query are often concentrated in specific segments rather than being uniformly distributed across the video. 
This presents a new challenge in allocating the total selection size \(k\) across multiple segments.
%
Focusing on video summarization, \citet{gong2014diverse} allocate \(k\) uniformly across segments,
while \citet{zheng2021k} formulate it as an integer programming problem and relax the constraints to find an approximate solution.
%
Instead, we focus on general VidQA and model this sequential DPP across the segment sequence as an MDP. 
We propose a dynamic programming algorithm that can optimally allocates the total selection size \(k\) across segments with pseudo-polynomial complexity, without relaxing problem constraints.


First, we segment a video into multiple segments of equal length \(m\), so that a video with \(n\) frames is divided into \(T = \lceil n/m \rceil\) segments. 
The corresponding frame indices of the video, \(\N_n\), are divided as \(\N_n = \bigcup_{t=1}^{T} N_t\), 
where \(N_t\) denotes the frame indices of the \(t\)-th segment, with $|N_t|=m$.
Moreover, let the selected subset from the \(t\)-th segment be denoted as \(S_t \subseteq N_t\).
%
Therefore, when applying DPP to the \(t\)-th segment, conditioned on the selection in the \((t-1)\)-th segment \(S_{t-1}\), 
the DPP conditional distribution~\cite{gong2014diverse} for the current segment \(S_t\) is expressed as:
\begin{equation}\label{eq:cond_dpp}
    \mathcal{P}(S_t | S_{t-1}) = \frac{\det(\tilde{\mL}_{S_{t-1} \cup S_t})}{\det(\tilde{\mL}_t + \I_t)} \,,
\end{equation}
where \(\tilde{\mL}_t\) is the submatrix of \(\tilde{\mL}\) for the indices \(S_{t-1} \cup N_t\), and \(\I_t\) is a diagonal matrix with zeros for \(S_{t-1}\) and ones for \(N_t\).
%
Thus, the joint distribution of subsets can be written as:
\begin{equation}\label{eq:seq_score}
    \mathcal{P}(S_1, S_2, \dots, S_T) 
        = \mathcal{P}(S_1) \prod_{t=2}^{T} \mathcal{P}(S_t \mid S_{t-1}) \,.
\end{equation}
%
If the selection size for the \(t\)-th segment is predetermined as \(|S_t| = k_t\),
the MAP inference in \cref{eq:seq_score} can be performed sequentially, from the first to the last~\cite{gong2014diverse}, similar to Bayesian belief updates (such as Kalman filtering):
\begin{equation}\label{eq:map_inf}
    S_t^* = \mathop{\mathrm{argmax}}_{S_t \subseteq N_t, |S_t|=k_t} \mathcal{P}(S_t \mid S_{t-1}^*) \,.
\end{equation}
This implies that when determining one of \(k_t\) and \(S_t^*\), the other is uniquely determined, \ie, \(k_t \leftrightarrows S_t^*\).


Determining the optimal allocation of selection size \(k_t\) for each segment is challenging. 
%
We observe that the MAP inference~(\cref{eq:map_inf}) of the DPP conditional distribution~(\cref{eq:cond_dpp}) uniquely depends on a triplet state, consisting of:
1)~the candidate indices in the current segment \(N_t\),
2)~the selections from the preceding segment \(S_{t-1}^*\), determined by $k_{t-1}$, and
3)~the remaining capacity for further selections \(k - C_{t-1}\), where \(C_{t-1}\) denotes the total selection size before the \((t-1)\)-th~(inclusive) segment.
Consequently, this MAP inference can be formulated as an MDP, where the state is defined as the triplet \((N_t, k_{t-1}, C_{t-1})\).
Moreover, the segment index \(t-1\) can uniquely determine \(N_t\).
Therefore, the state can be represented as \((t \in [1, T], k_t \in [0, k], C_t \in [0, k])\).
The state transition from \(t-1\) to \(t\) is deterministic when the action of selecting \(k_t\) frames from the \(t\)-th segment is given, which can be expressed as:
\begin{equation}
	(t-1, k_{t-1}, C_{t-1})
	\xrightarrow{\hspace{-0.1cm}
		\scriptsize\begin{array}{c}
	        \text{selecting } S_t \\
	        \text{with } |S_t|=k_t
    	\end{array}
	\hspace{-0.1cm}}
	(t, k_t,\overbrace{C_{t-1}+ k_t}^{C_t})\,.
\end{equation}
The corresponding reward in MDP is defined as
\begin{equation}\label{eq:reward}
    \R_{S_t}(t-1, k_{t-1}, C_{t-1}) = \log \left(\mathcal{P}(S_t \mid S_{t-1}^*)\right) \,,
\end{equation}
where \(S_{t-1}^*\) is uniquely determined by \(k_{t-1}\) in the state $(t-1, k_{t-1}, C_{t-1})$.
This reward can be efficiently computed using a greedy DPP algorithm applied to the conditional DPP (\ie, \cref{eq:cond_dpp}).
%
Therefore, the list-wise \(\mathrm{Score}\) function in \cref{eq:score}, based on \cref{eq:seq_score}, is defined as:
\begin{equation}\label{eq:opt_score}
	\mathrm{Score}(S_1,S_2,\dots,S_T) =
		\sum_{t=0}^{T-1} \R_{S_{t+1}}(t, k_t, C_t)\,.
\end{equation}
Using this definition and the \((1 - 1/e)\)-approximate greedy DPP algorithm, we can apply a dynamic programming approach to identify a solution \(\hat{S}\) such that \mbox{\(\mathrm{Score}(\hat{S}) \geq (1 - 1/e) \cdot \mathrm{Score}(S^*)\)}, with a closely \(\mathcal{O}(nk^4)\) time complexity~(proof provided in the Appendix).

We define a value function \(Q(t, k_t, C_t)\) as the maximum accumulated reward from the initial state \((0, 0, 0)\) to the current state \((t, k_t, C_t)\). 
Consequently, 
\begin{equation}
    Q^*(t, C_t) = \max\limits_{0 \leq k_t \leq k - C_t} Q(t, k_t, C_t)
\end{equation}
represents the maximum accumulated reward for selecting \(C_t\) frames up to and including the \(t\)-th segment. 
The solution to \cref{eq:opt_score} is given by \(\hat{S} = \mathop{\mathrm{argmax}}_{S} Q^*(T, k)\).

\begin{algorithm}[!t]
    \caption{\label{alg} \mname}
    \KwIn{Video $V = \{f_i\}_{i=1}^{n}$, query $q$, select size $k$.}
    \KwOut{A set of frame indices $S$ with $|S| = k$.}
    
    Initialize $Q_{0,0}^* = 0$ and $\T_{0,0} = [\,]$\;
    Extract embeddings $f_i$ and $q$ using \cref{eq:emb}\;
    
    \For{$t \gets 1$ to $T = \lceil n/m \rceil$}{
        Compute $\mr_t$, $\mL_t$, and $\tilde{\mL}_t$ using \cref{eq:k_family,eq:cond_k,eq:tilde_L,eq:cmgk}\;
        Compute offset $\mathrm{o} = -\log(\det(\tilde{\mL}_t + \I_t))$\;
        
        \For{$C_{t-1} \gets 0$ to $k$}{
            \For{$k_{t-1} \gets 0$ to $\min(m, k, k - C_{t-1})$}{
            	Get $C_t = C_{t-1}+k_{t-1}$\;
	            Get $\mathrm{rwd}^* = \R_{S_t}^*(t-1, C_{t-1})$ and $S_t = \pi(t, C_t)$ using DPP as \cref{eq:greedy,eq:cond_dpp,eq:seq_score,eq:reward}\;
                Compute $cur\_q = Q_{t-1, C_{t-1}} + \mathrm{rwd}^*$\;

                \If{$cur\_q > Q_{t, C_t}$}{
                    Update $Q_{t, C_t} \gets cur\_q$\;
                    Update $\T_{t, C_t} \gets \T_{t-1, C_{t-1}} \cup S_t$\;
                }
            }
        }
    }
    \Return $\T_{T, k}$\;
\end{algorithm}



\begin{table*}[!t]
\centering
\begin{tabular}{lcccccccc}
    \toprule
		\multirow{2}{*}{\textbf{Model}} & 
		\multirow{2}{*}{\makecell{\textbf{LLM} \\ \textbf{Size}}} & 
		\multirow{2}{*}{\textbf{\#Frames}} & 
		\multicolumn{4}{c}{\textbf{Video-MME$_{(wo/w\ subs)}$}} & 
		\multirow{2}{*}{\textbf{MLVU}} &
		\multirow{2}{*}{\textbf{LVB$_{val}$}} \\

		\cmidrule(lr){4-7} & & & Overall & Short & Medium & Long & & \\ 
		
		\multicolumn{2}{c}{\textit{Avg. Video Duration}} && \textit{17min} & \textit{1.3min} & \textit{9min} & \textit{41min} & \textit{12min} & \textit{12min} \\
  	\toprule
	Video-LLaVA     & 7B & 8 & 39.9 / 41.6 & 45.3 / 46.1 & 38.0 / 40.7 & 36.2 / 38.1 & 47.3 &   - \\
	Qwen-VL-Chat    & 7B & 8 & 41.1 / 41.9 & 46.9 / 47.3 & 38.7 / 40.4 & 37.8 / 37.9 &   -  &   -  \\
	VideoChat2      & 7B & 8 & 39.5 / 43.8 & 48.3 / 52.8 & 37.0 / 39.4 & 33.2 / 39.2 & 44.5 &   -  \\
	Chat-UniVi-V1.5 & 7B & 8 & 40.6 / 45.9 & 45.7 / 51.2 & 40.3 / 44.6 & 35.8 / 41.8 &   -  &   -  \\
	VideoLLaMA2     & 7B & 8 & 47.9 / \unk & 56.0 / \unk & 45.4 / \unk & 42.1 / \unk &   -  &   -  \\
	LLaVA-NeXT-QW2  & 7B & 8 & 49.5 / \unk & 58.0 / \unk & 47.0 / \unk & 43.4 / \unk &   -  &   -  \\
	\hdashline
	LongVILA        & 8B & 128 & 49.2 / \unk & 60.2 / \unk & 48.2 / \unk & 38.8 / \unk &   -  &   -  \\
	LongVA          & 7B & 128 & 52.6 / 54.3 & 61.1 / 61.6 & 50.4 / 53.6 & 46.2 / 47.6 &   -  &   -  \\ 
	Video-XL        & 7B & 128/256 & 55.5 / 61.0 & 64.0 / 67.4 & 53.2 / 60.7 & 49.2 / 54.9 & 64.9 &   -  \\
	LLaVA-OneVision$^*$ 
					& 7B & 32 & 58.2 / \unk & \unk / \unk & \unk / \unk & \unk / \unk & 64.7 & 56.3 \\
	\midrule
	VILA-V1.5       & 8B & 8 & 47.5 / 50.0 & 57.8 / 61.6 & 44.3	/ 46.2 & 40.3 / 42.1 & 46.3 & 47.1 \\
	\rule{0pt}{2ex} +\textsc{Frame-Voyager}\xspace
				    & 8B & 8 & 50.5 / 53.6 & 60.3 / 65.0 & 47.3 / 50.3 & 43.9 / 45.3 & 49.8 &   -  \\
	\rowcolor{gray!30}   
	\rule{0pt}{2ex} \textbf{+\mname}
	                & 8B & 8 & 53.3	/ 56.6 & 65.7 / 68.7 & 49.9 / 54.4 & 44.2 / 46.6 & 58.6 & 50.8 \\
	\hdashline
	MiniCPM-V2.6    & 7B & 8 & 52.6	/ 53.1 & 61.1 / 63.8 & 50.3	/ 50.2 & 46.4 / 45.4 & 55.4 & 51.2 \\
	\rowcolor{gray!30}
	\rule{0pt}{2ex} \textbf{+\mname} 
	                & 7B & 8 & 58.0 / 61.8 & 69.1 / 71.6 & 56.4 / 61.7 & 48.4 / 52.1 & 66.6 & 57.1 \\
	\hdashline
	LLaVA-OneVision & 7B & 8 & 53.6 / 53.9 & 64.7 / 65.0 & 51.4 / 52.0 & 44.7 / 44.6 & 59.3 & 54.2 \\
	\rule{0pt}{2ex} +\textsc{Frame-Voyager}\xspace
				    & 7B & 8 & 57.5 / \unk & 67.3 / \unk & 56.3 / \unk & 48.9 / \unk & 65.6 &   - \\
	\rowcolor{gray!30}	
	\rule{0pt}{2ex} \textbf{+\mname} 
	                & 7B & 8 & 59.6 / 59.1 & 71.8 / 71.1 & 57.6 / 57.6 & 49.6 / 48.8 & 69.8 & 59.0 \\
	\hdashline
	Ovis2 & 7B & 8 & 58.9 / 62.1 & 67.8 / 72.0 & 57.4 /60.3 & 51.3 / 54.0 & 60.9 & 56.9 \\
	\rowcolor{gray!30}
	\rule{0pt}{2ex} \textbf{+\mname} 
	                & 7B & 8 & \textbf{63.9 / 66.1} & \textbf{75.3 / 77.6} & \textbf{63.2 / 64.3} & \textbf{53.2 / 56.6} & \textbf{73.9} & \textbf{62.7} \\
	\bottomrule
	\end{tabular}
	\caption{
%	    Comparison of Video-LLMs with and without \mname{} as a plug-and-play process. 
%    	Results are shown for the Video-MME benchmark under two settings: without subtitles (w/o sub.) and with subtitles (w sub.). 
%	    LVB$_{val}$ denotes the LongVideoBench validation set, excluding interleaved subtitles. 
%	    %
%	    Video-XL uses 128 frames for Video-MME and 256 frames for MLVU.
%	    LLaVA-OneVision$^*$ refers to results from the official report with the well-tuned number of frames.
%	    The accuracy percentage sign ($\%$) has been omitted for clarity.
%		Comparison of Video-LLMs with and without \mname{}.
%    	Results are shown for the Video-MME under two settings: without subtitles and with subtitles. 
%	    LVB$_{val}$ denotes the LongVideoBench validation set, excluding interleaved subtitles. 
%	    %
%	    Video-XL uses 128 frames for Video-MME and 256 frames for MLVU.
%	    LLaVA-OneVision$^*$ refers to results from the official report with the well-tuned frame counts.
	Comparison of Video-LLMs with and without \mname{}.
	LVB$_{val}$ denotes the LongVideoBench validation set without subtitle interleaving.
	Video-XL uses 128/256 frames on Video-MME/MLVU.
	LLaVA-OneVision$^*$ refers to official report with well-tuned frames.
	}
	\label{tab:main_results}
\end{table*}


Dynamic programming can be used to update the \(Q\) score and identify the solution \(\hat{S}\) with an optimal allocation of \(k\), with a time complexity closely \(\mathcal{O}(nk^4)\).
%
In practice, to improve efficiency further, we adopt lazy dynamic programming strategies.
First, we directly maintain the \(Q^*\) score instead of updating \(Q\).
This strategy assumes that selecting a frame from the \(t\)-th segment depends only on the preceding state with optimal $k_{t-1}^*$, \ie, \((t-1, k_{t-1}^*, C_{t-1})\), corresponding to \(Q^*(t-1, C_{t-1})\), similar to the online MAP inference described in \cref{eq:map_inf}~\cite{gong2014diverse}.
Second, only the most recently selected frame is used as the condition in~\cref{eq:reward}, approximating an adaptive condition size and a fixed candidate segment size \(m\) for the \(t\)-th segment.
Consequently, with the lazy strategies, the state is simplified to \((t, C_t)\).
The policy function is denoted by \(S_t^*=\pi(t, C_t)\), and \(\T_{t, C_t}\) is an initially empty list \([\,]\) that tracks the selection trace \([S_1, S_2, \dots, S_t]\).
%
The corresponding reward simplifies to:
\begin{equation}
	\R_{S_t}^*(t-1,C_t-k_t)=\log (\mathcal{P}(S_t \mid \T_{t-1, C_{t-1}}[-1])), 	
\end{equation}
where $C_t-k_t = C_{t-1}$.
For brevity, we denote $Q^*_{t, C_t}=Q^*(t, C_t)$.
The dynamic programming update is:
\begin{align}\label{eq:dp_update}
	&Q^*_{t, C_t} = \max_{S_t \subseteq N_t} Q^*_{t-1, C_t-k_t} + \R_{S_t}^*(t-1,C_t-k_t)\,, \notag \\
	&\pi(t, C_t)  = \mathop{\mathrm{argmax}}_{S_t \subseteq N_t} 
					Q^*_{t-1, C_t-k_t} + \R_{S_t}^*(t-1,C_t-k_t)\,, \notag \\
	&\T_{t, C_t} = \T_{t, C_t - |S_t^*|} + [S_t^*],\text{where } S_t^* = \pi(t, C_t)\,.
\end{align}
The final solution is \(\hat{S} = \mathop{\mathrm{argmax}}_{S} Q^*(T, k) = \T(T, k) \).


The pseudocode for \mname{} is provided in \cref{alg}.
The end index \((t, C_t + k_t)\) is iterated, while the start index \((t-1, C_t)\) remains fixed.
The update functions are equivalent with \cref{eq:dp_update}.
This implementation is more efficient as the greedy DPP algorithm (line 9) can be computed incrementally in $\mathcal{O}(mk)$.
Thus, the worst-case time complexity of \cref{alg} is closely \(\mathcal{O}(nk^3)\) in practice~(proof in Appendix).
Compared to the standard DPP complexity of \(\mathcal{O}(nk^2)\), the additional factor \(k \ll n\) is negligible. 
Additionally, the iteration in line 6 is independent, enabling parallel updates.
This results in a time complexity of \(\mathcal{O}(nk^2)\), making its efficiency comparable to the standard DPP.




